\section{Chemistry Reaction
Study\label{ch5:chemistryReactionStudy}}

\todo{Use Results from AIAAA, laminar comparison between frozen and
nonequilibrium flow}
This section investigates nonequilibrium effects of aero-optics in a laminar
flow by comparing a flow with nonequilibrium chemistry and the other one with
frozen chemistry (chemical composition of the gas remains unchanged because
the reactions are too slow to occur within the times scale of the flow). \TabRef{tab:chemistryReactionStudy} shows the cases used for this
investigation.  

\begin{table}[H]
    \begin{center}
        \begin{tabular}{|c|c|c|c|c|c|c|}
            \hline
            Case & $M\Units{\;}$ & $T_{wall}\Units{K}$ &
            $T_{\infty}\Units{K}$ & $P_{\infty}\Units{Pa}$ &
            $h_0\Units{MJ}$ & $q\Units{kPa}$ \\
            \hline\hline
            $1C$\cite{Mackey2019} & $11.2$ & $1000$ & $293$ & $1200$ & $7.38$  & $105$ \\\hline
            $2C$\cite{Mackey2019} & $13.2$ & $1000$ & $293$ & $1200$ & $10.25$ & $146$ \\\hline
            $3C$\cite{Mackey2019} & $15.2$ & $1000$ & $293$ & $1200$ & $13.60$ & $193$ \\\hline
        \end{tabular}
    \caption{Chemistry Study \label{tab:chemistryReactionStudy}}
    \end{center}
\end{table}

    For this investigation the IRV-2 vehicle, shown in \FigRef{fig:IRV2} is used.
    The IRV-2 vehicle is a sphere-biconic with a nose radius of $1.95\Units{cm}$ and a total
    length of $1.386 \Units{m}$. Six drag flaps are located at the aft end of
    the vehicle and can be deployed to a maximum angle of $45$ degrees
    \cite{Kuntz2001}. For this work, only axisymmetric calculations on the
    IRV-2 nosetip are performed. \FigRef{fig:IRV2Mesh} shows
    the mesh used in this work. 

    \begin{figure}[H]
        \centering
        \begin{minipage}{0.49\textwidth}
            \centering
            \includegraphics[scale=0.5]{chapter5/chemistryReaction/irv2}
            \caption{IRV-2 vehicle}\label{fig:IRV2}
        \end{minipage}
        \hfill
        \begin{minipage}{0.49\textwidth}
            \centering
            \includegraphics[scale=0.3]{chapter5/chemistryReaction/1C_frozenMesh}
            \caption{IRV-2 mesh}\label{fig:IRV2Mesh}
        \end{minipage}
    \end{figure}

    \newpage % Translational Temperatures
        \begin{landscape}
            \begin{figure}[H]
                \centering
                \begin{subfigure}[Case $1C$]
                {
                    \label{fig:frozenTr1C}
                    \includegraphics[width=0.4\textwidth]{chapter5/chemistryReaction/1C_frozenContourTtr.png}
                }
                \end{subfigure}
                \hfill
                \begin{subfigure}[Case $2C$]
                {
                    \label{fig:frozenTr2C}
                    \includegraphics[width=0.4\textwidth]{chapter5/chemistryReaction/2C_frozenContourTtr.png}
                }
                \end{subfigure}
                \hfill
                \begin{subfigure}[Case $3C$]
                {
                    \label{fig:frozenTr3C}
                    \includegraphics[width=0.4\textwidth]{chapter5/chemistryReaction/3C_frozenContourTtr.png}
                }
                \end{subfigure}
                \caption{Contour translational temperature for
                frozen cases from \TabRef{tab:chemistryReactionStudy}}
            \end{figure}

            \begin{figure}[H]
                \centering
                \begin{subfigure}[Case $1C$]
                {
                    \label{fig:nonequilibriumTtr1C}
                    \includegraphics[width=0.4\textwidth]{chapter5/chemistryReaction/1C_nonequilibriumContourTtr.png}
                }
                \end{subfigure}
                \hfill
                \begin{subfigure}[Case $2C$]
                {
                    \label{fig:nonequilibriumTtr2C}
                    \includegraphics[width=0.4\textwidth]{chapter5/chemistryReaction/2C_nonequilibriumContourTtr.png}
                }
                \end{subfigure}
                \hfill
                \begin{subfigure}[Case $3C$]
                {
                    \label{fig:nonequilibriumTtr3C}
                    \includegraphics[width=0.4\textwidth]{chapter5/chemistryReaction/3C_nonequilibriumContourTtr.png}
                }
                \end{subfigure}
                \caption{Contour translational temperature for
                nonequilibrium cases from \TabRef{tab:chemistryReactionStudy}}
            \end{figure}
        \end{landscape}
        
    \newpage % Vibrational Temperatures
        \begin{landscape}
            \begin{figure}[H]
                \centering
                \begin{subfigure}[Case $1C$]
                {
                    \label{fig:frozenTvr1C}
                    \includegraphics[width=0.4\textwidth]{chapter5/chemistryReaction/1C_frozenContourTvr.png}
                }
                \end{subfigure}
                \hfill
                \begin{subfigure}[Case $2C$]
                {
                    \label{fig:frozenTvr2C}
                    \includegraphics[width=0.4\textwidth]{chapter5/chemistryReaction/2C_frozenContourTvr.png}
                }
                \end{subfigure}
                \hfill
                \begin{subfigure}[Case $3C$]
                {
                    \label{fig:frozenTvr3C}
                    \includegraphics[width=0.4\textwidth]{chapter5/chemistryReaction/3C_frozenContourTvr.png}
                }
                \end{subfigure}
                \caption{Contour vibrational temperature for
                frozen cases from \TabRef{tab:chemistryReactionStudy}}
            \end{figure}

            \begin{figure}[H]
                \centering
                \begin{subfigure}[Case $1C$]
                {
                    \label{fig:nonequilibriumTvr1C}
                    \includegraphics[width=0.4\textwidth]{chapter5/chemistryReaction/1C_nonequilibriumContourTvr.png}
                }
                \end{subfigure}
                \hfill
                \begin{subfigure}[Case $2C$]
                {
                    \label{fig:nonequilibriumTvr2C}
                    \includegraphics[width=0.4\textwidth]{chapter5/chemistryReaction/2C_nonequilibriumContourTvr.png}
                }
                \end{subfigure}
                \hfill
                \begin{subfigure}[Case $3C$]
                {
                    \label{fig:nonequilibriumTvr3C}
                    \includegraphics[width=0.4\textwidth]{chapter5/chemistryReaction/3C_nonequilibriumContourTvr.png}
                }
                \end{subfigure}
                \caption{Contour vibrational temperature for
                nonequilibrium cases from \TabRef{tab:chemistryReactionStudy}}
            \end{figure}
        \end{landscape}
        
    \subsection{Results Case $1C$}
    \subsection{Results Case $2C$}
    \subsection{Results Case $3C$}


